\documentclass[a4paper]{article}
\usepackage{amsmath}
\usepackage{amssymb}
\usepackage{geometry}
\usepackage{hyperref}
\usepackage{booktabs}
\usepackage{enumerate}
\usepackage{listings}
\usepackage{xcolor}
\geometry{margin=1in}

\lstdefinestyle{python}{
    basicstyle=\ttfamily\small,
    keywordstyle=\color{blue},
    commentstyle=\color{gray},
    stringstyle=\color{red},
    frame=single,
    breaklines=true,
    showstringspaces=false,
}

\begin{document}

\title{\textit{Introduction to Econometrics}\\Assignment 6}
\author{Richard Harnisch, s5238366}
\date{\today}
\maketitle

% answer environment
\makeatletter
\newcommand{\answer}[1]{%
    \par\medskip
    \noindent
    \hspace*{-\@totalleftmargin}%
    \fbox{%
        \parbox{\dimexpr\textwidth-2\fboxsep-2\fboxrule}{#1}%
    }%
    \par\medskip
}
\makeatother

\noindent\textit{TeX code available under\\ \url{https://github.com/richardharnisch/intro-to-econometrics}}

\begin{enumerate}

    \item Consider the model $y_1 = \mathbf{x}_1\delta_1 + \alpha_1 y_2 + \gamma q + \varepsilon_1$
          where $q$ is omitted. Assume $E(\varepsilon_1 \mid \mathbf{x}_1, y_2, q) = 0$,
          $E(q \mid \mathbf{x}_1) = 0$, but $y_2$ and $q$ are correlated.
          Let $\mathbf{x}_2$ be instruments for $y_2$ with relevance and
          $E(\varepsilon_1 \mid \mathbf{x}_2) = 0$, yet $\mathbf{x}_2$ may be
          correlated with $q$.

          Would OLS and IV be consistent if $y_2$ and $\mathbf{x}_2$ are
          correlated with $q$?

          \answer{
              Neither OLS nor IV would be consistent.

              For OLS, the estimated model is $y_1 = \mathbf{x}_1\delta_1 + \alpha_1 y_2 + u$
              where $u = \gamma q + \varepsilon_1$ (the error absorbs the omitted $q$). Since
              $\text{Cov}(y_2, q) \neq 0$ and $\gamma \neq 0$, we have
              $\text{Cov}(y_2, u) = \gamma\text{Cov}(y_2, q) \neq 0$,
              so $\hat{\alpha}_1^{\text{OLS}} \xrightarrow{p} \alpha_1 + \text{bias} \neq \alpha_1$.

              For IV to be consistent we require $E(\mathbf{x}_2 u) = 0$, i.e.
              \[
                  E(\mathbf{x}_2 [\gamma q + \varepsilon_1]) = 0.
              \]
              We are given $E(\mathbf{x}_2 \varepsilon_1) = 0$, but if $\mathbf{x}_2$ is
              correlated with $q$ then $E(\mathbf{x}_2 q) \neq 0$, which means
              $E(\mathbf{x}_2 \gamma q) \neq 0$. The exclusion restriction is therefore
              violated and IV is also inconsistent. The root cause is that $q$ is omitted
              and $\mathbf{x}_2$ is not orthogonal to it; controlling for $q$ or finding
              instruments that satisfy the exclusion restriction with respect to $q$ would
              be required for consistency.
          }

    \item We estimate returns to education using the AK1991 dataset.

          \begin{enumerate}

              \item (1 point) Write the model to estimate returns to education.

                    \answer{
                        \[
                            \log(\text{wage}_i) = \beta_0 + \beta_1 \text{education}_i
                            + \beta_2 \text{married}_i + \beta_3 \text{black}_i
                            + \beta_4 \text{smsa}_i
                            + \sum_{s=2}^{50} \gamma_s \text{state}_{s,i}
                            + \sum_{y=1931}^{1939} \delta_y \text{yob}_{y,i} + \varepsilon_i
                        \]
                        Here $\text{wage}_i$ denotes weekly wages, $\text{education}_i$ years of
                        education, $\text{married}_i$ and $\text{black}_i$ are binary indicators,
                        and $\text{smsa}_i$ indicates residence in a Standard Metropolitan
                        Statistical Area. The model includes a full set of state indicators
                        (49 dummies, state 1 omitted) and a full set of year-of-birth indicators
                        (9 dummies, 1930 omitted). The parameter of interest is $\beta_1$,
                        the return to one additional year of education.
                    }

              \item (1 point) What is the potential issue with OLS?

                    \answer{
                        The main concern is endogeneity caused by omitted ability. Unobserved
                        ability (intelligence, motivation, etc.) affects both how much education
                        an individual acquires and how much they earn, so it is absorbed into
                        the error term $\varepsilon_i$. Because higher-ability individuals tend
                        to acquire more education and earn higher wages,
                        $\text{Cov}(\text{education}_i, \varepsilon_i) > 0$. This means
                        $\hat{\beta}_1^{\text{OLS}}$ is biased upward: it conflates the causal
                        effect of education with the effect of ability, and is therefore not a
                        consistent estimator of the return to education.
                    }

              \item (2 points) Is education related to quarter of birth?
                    Could QOB be used as instrument?

                    \answer{
                        Yes, education is significantly related to quarter of birth. The
                        first-stage regression regresses years of education on the QOB dummies
                        and all controls from the main equation (married, black, smsa, full
                        state dummies, full year-of-birth dummies). The results are:

                        \begin{center}
                            \begin{tabular}{lccc}
                                \toprule
                                QOB & Coefficient & SE & t-statistic \\
                                \midrule
                                QOB=2 & 0.0457 & 0.0157 & 2.913 \\
                                QOB=3 & 0.0982 & 0.0154 & 6.390 \\
                                QOB=4 & 0.1454 & 0.0156 & 9.295 \\
                                \midrule
                                \multicolumn{4}{c}{F-test for joint significance: $F_{(3,329443)} = 32.65$, $p < 0.001$} \\
                                \bottomrule
                            \end{tabular}
                        \end{center}

                        The pattern is explained by compulsory schooling laws, which require
                        students to remain in school until age 16 or 17. Since school entry
                        occurs in September, children born in the first quarter are the youngest
                        in their cohort. When they reach the compulsory leaving age they have
                        completed fewer years of schooling than children born later in the year,
                        creating exogenous variation in education by quarter of birth.

                        Quarter of birth is a valid instrument. The relevance condition is
                        satisfied: the first-stage F-statistic of 32.65 far exceeds the
                        conventional threshold of 10, indicating a strong relationship between
                        QOB and education. The exogeneity condition is also plausibly satisfied:
                        season of birth is determined at conception, beyond individual control,
                        and is unlikely to be related to ability, family background, or other
                        determinants of wages except through its effect on years of schooling.
                    }

              \item (3 points) Perform a fully-fledged IV analysis with tests.

                    \answer{
                        The table below compares OLS and 2SLS estimates of the return to
                        education, using QOB dummies as instruments.

                        \begin{center}
                            \begin{tabular}{lrr}
                                \toprule
                                & OLS & 2SLS \\
                                \midrule
                                Education (return/year) & 0.0625 & 0.0968 \\
                                Standard Error & 0.0003 & 0.0209 \\
                                t-statistic & 182.06 & 4.64 \\
                                p-value & $<$0.0001 & $<$0.0001 \\
                                \midrule
                                \multicolumn{3}{c}{95\% CI for 2SLS: $[0.0559,\ 0.1378]$} \\
                                \bottomrule
                            \end{tabular}
                        \end{center}

                        According to 2SLS, one additional year of education increases weekly
                        wages by approximately 9.68\%, compared to 6.25\% under OLS. The
                        higher IV estimate is consistent with classical measurement error in
                        education attenuating the OLS coefficient toward zero, or with the IV
                        recovering a local average treatment effect for compliers whose
                        education is affected by compulsory schooling laws and who may face
                        higher marginal returns than average.

                        Three diagnostic tests were performed. First, the first-stage
                        F-statistic of 32.65 comfortably exceeds 10, so there is no weak
                        instrument problem. Second, the Hausman test for endogeneity
                        ($H_0$: education is exogenous) yields $t = 1.643$ with $p = 0.100$,
                        so we fail to reject exogeneity at the 5\% level; the difference
                        between OLS and 2SLS is not statistically significant, though this
                        does not rule out endogeneity. Third, the Sargan test for
                        overidentifying restrictions uses the fact that with 3 instruments
                        for 1 endogenous variable there are $L - K = 3 - 1 = 2$
                        overidentifying restrictions. The test statistic is 2.244,
                        well below the 5\% critical value $\chi^2_{2,0.05} = 5.991$
                        ($p = 0.523$), so we cannot reject the hypothesis that all
                        exclusion restrictions hold.

                        In conclusion, the IV estimate of 9.68\% return per year of education
                        is statistically significant, rests on a strong first stage, and is
                        supported by the Sargan test. Although the Hausman test does not
                        conclusively establish endogeneity, the IV approach is motivated by
                        the institutional setting in which compulsory schooling laws generate
                        plausibly exogenous variation in years of schooling.
                    }

          \end{enumerate}

\end{enumerate}

\end{document}
